
% =====================================================================
\section{Data, institutions, and the sample}
\label{sec:data}
% =====================================================================

\subsection{The feed}

The source is the Nikkei NEEDS \texttt{TICST120} product for calendar 2024:
one record per market event, carrying the trade if there was one and the ten best
price levels on each side of the book as they stood afterwards. Trade direction
is supplied by the exchange rather than inferred, which removes a source of error
that would otherwise sit underneath every result here. Daily summaries from the
companion \texttt{TICSS110} product supply trading units, shares outstanding,
and the variables used to screen the universe.

Five dates are excluded at the calendar level. Four never arrived in the
delivered feed, and 2024-04-23 arrived with ten shards against a monthly median
of twenty-two --- the signature of a transfer that stopped part-way. Treating a
partial day as a quiet one would bias every measure computed from it, so the
exclusion is made once, at the calendar, and no later stage can rediscover them.
That leaves 240 usable trading days.

\subsection{Two institutional details that shape everything}

\paragraph{The session changed length part-way through the year.}
The afternoon session closed at 15:00 until 2024-11-05 and at 15:30 afterwards.
Every time-dependent quantity here is parameterised by date. A hardcoded close
would silently truncate two months of the afternoon, and would corrupt the
price-impact horizons on exactly those days.

\paragraph{The exchange runs two tick grids at once.}
Constituents of the TOPIX 500 trade on a finer grid than everything else --- a
regime extended from the TOPIX 100 to the whole index on 2023-06-05 and stable
through 2024. The schedule is encoded from the exchange's published tables and
cross-checked against ticks read directly off the tape.

This interacts with Ohta's sample filter in a way worth stating plainly, because
it shapes the sample more than any other decision. The filter admits only tick
sizes that are powers of ten, since the last-digit-of-ten construction is
otherwise undefined --- on a 0.5-yen grid the last digit takes only even values
and ``round'' stops meaning what it means elsewhere. On the finer grid the
1,000--3,000 yen band is 0.5 yen, so a large part of the mid-priced TOPIX 500
population is excluded outright. The regression sample is therefore tilted away
from mid-priced large caps, and the composition is reported rather than papered
over:

\begin{table}[htbp]
\centering
\caption{Tick-size composition of evaluated stock-days}
\label{tab:tickcomp}
\small
\begin{tabular}{@{}lrl@{}}
\toprule
Tick size & Stock-days & Status \\
\midrule
1 yen & 686 & kept \\
5 yen & 33 & excluded by filter (d) \\
10 yen & 28 & kept \\
0.1 yen & 26 & kept \\
0.5 yen & 24 & excluded by filter (d) \\
\bottomrule
\end{tabular}
\end{table}

\subsection{Building the sample}

The universe is selected by Ohta's own criteria, applied first to the daily
summaries so that the expensive tick-level work runs only on stocks that could
enter the sample. He studies First Section (now Prime) common stocks whose
trading days satisfy four conditions on more than half the year's sessions: the
first trade by 9:10, an opening price above 200 yen, more than twenty
continuous-session trades, and a tick size that is a power of ten throughout the
day. Regional venues, exchange-traded funds and trust units --- identifiable by
their trading units --- are excluded.

Each of those filters exists for a reason worth keeping in view. The 9:10 rule
removes days that opened late after a price-limit halt, which would have a
shortened continuous session and therefore a differently-behaved measure. The
200-yen floor removes stocks so cheap that the price cannot traverse all ten
digits without moving several percent. The twenty-trade minimum removes days
where the measure is a ratio of small integers.

Applying the same tests properly at tick level gives the final panel:
800 stock-days attempted, 692
passing all four filters, and 668 surviving the requirement that a stock
qualify on more than half the year --- 334 stocks over 2 trading
days, from 2024-01-04 to 2024-01-05. The full waterfall is reported in the stage output; no
stock-day disappears without being counted.

\paragraph{Verdict.}
The sample is Ohta's, reconstructed on a year he did not have. Its one
distinctive feature --- the exclusion of the 0.5-yen grid and therefore of much
of the mid-priced large-cap population --- follows from his filter rather than
from a choice made here, and is carried explicitly into the interpretation of
every result below.
